\begin{afterword}
\summaryhead

The post names underconfidence as the third of three besetting sins of rationalists. People warned about a bias sometimes correct too far, and someone who has read many studies of overconfidence may push their self-estimate too low, especially since modesty feels safer and kinder than pride. The danger is passing up what one could have done.

The first advice is to test hypotheses about one's own abilities, as the author did with the AI-Box Experiment, winning once and losing twice when the stakes were high. Smart people who are used to winning avoid hard problems; one should try to win every time, yet seek challenges hard enough to lose sometimes. The author's example is a proposed debate, without preparation, with a theist said to have beaten Christopher Hitchens.

A subtler form is lost momentum: doubts about oneself raised and never tested. The last advice is to ask of any habit of thought, ``Does this way of thinking make me stronger, or weaker?'', with signs to watch for. The author learned this from a childhood math test, where second-guessing lowered the scores. Underconfidence, the post ends, is a ``stopping mistake'' that prevents the experience that would correct it.

\responsehead

The post is advice drawn from the author's experience, and it is careful with its claims. Its premise, that people warned about a bias sometimes overcorrect, is hedged (``not unheard of'') and supported by research on bias correction, which reports that corrections ``often'' produce a bias in the other direction. Its main remedy is sound and concrete: treat a doubt about your ability as a hypothesis and test it. The last question, ``stronger, or weaker?'', could have been a matter of feeling, but the post attaches signs one can observe, such as failing less often or avoiding challenges.

The author's own cases are accurately reported. The AI-Box record matches the account in ``Shut Up and Do the Impossible!'': one win and two losses in the high-stakes games. The explanation from ``intelligence'' and ``effort'' is offered only as ``It is said''; it comes from Carol Dweck's research, which the notes describe. The math-test story runs against the general finding that changed answers usually improve scores, but the post draws only a hedged lesson from it and says the test was needed, which is its own argument for testing.

The debate is the one case in which the author changed course. The theist was William Lane Craig. Commenters urged the author to prepare, and the next day the author wrote that the news that Richard Carrier had also lost to Craig was ``the main thing that moved me,'' since the debate ``might still be a test of my ability even if I prepare.'' Craig later declined, so the test never took place.

Two points depend on the author's own view. The ranking of underconfidence among the sins is the author's impression, hedged (``does seem''), of the aspiring rationalists the author meets. And the post calls three ``reasonable-sounding'' positions ``impediments to the Way''; the first, two-boxing on Newcomb's problem, is still an open dispute among philosophers, and the post points to its earlier arguments rather than making them here.

The post does not name the first two sins. ``Final Words,'' placed earlier in the book, has Jeffreyssai name three flaws, the third of which is underconfidence.

In short: sound, testable advice against overcorrecting for bias, supported by research on overcorrection and by the author's reported cases, with its ranking of underconfidence offered as the author's impression.
\end{afterword}
